\chapter{Appendix: The Calculus and Linear Algebra You Need}
\label{chAcla:start}

Probability is built on two older subjects, and it leans on them constantly. Calculus is the
language of \emph{continuous} chance: the moment you write a probability density, an expectation
becomes an integral, a tail probability becomes an area, and the normal curve becomes a function you
actually have to integrate. Linear algebra is the language of \emph{many variables at once}: a vector
of random quantities, a covariance \emph{matrix}, a Markov chain's transition matrix, a principal
component, a least-squares prediction --- all of these are linear-algebra objects wearing
probabilistic clothes.

This appendix is not a course. It is a focused, intuition-first refresher on exactly the calculus and
linear algebra the rest of this book uses, and nothing more. Where a topic deserves a full treatment,
we point you to the companion volumes --- \emph{Calculus From the Ground Up} and \emph{Linear Algebra
From the Ground Up} --- which develop these ideas with the same care we bring to probability here. Our
goal is narrower: to make sure that when the chapter on expectation writes $\E[X]=\int x f(x)\,\dd x$,
or the chapter on the multivariate normal writes $\Sigma = Q \Lambda Q^\top$, you see exactly what
those symbols are doing and why they are allowed to.

We will move in two halves. First calculus: derivatives as rates and slopes, integrals as
accumulated area, the Fundamental Theorem that ties them together, multiple integrals over regions of
the plane, and gradients for functions of several variables. Then linear algebra: vectors and the dot
product, matrices as linear maps, eigenvectors and eigenvalues, the special role of symmetric and
positive-definite matrices, and orthogonal projection. Throughout, we keep one eye on probability and
flag the exact spot in the main text where each tool gets used.

\begin{focusbox}
You do not need to be fluent in calculus and linear algebra to read this book. You need to be
\emph{comfortable} with a specific short list of ideas. This appendix is that list, told as a story.
If a section feels obvious, skim it; if a section feels new, slow down --- it will pay off the next
time you meet an integral or a covariance matrix in the wild.
\end{focusbox}

\section{Part I: The Calculus You Need}
\label{chAcla:calc-intro}

\subsection{Functions, slopes, and the derivative}
\label{chAcla:derivative}

Start with the most basic object: a function $f$ that takes a number $x$ and returns a number $f(x)$.
Picture its graph, the curve $y=f(x)$. The first question calculus asks is: \emph{how fast is $f$
changing} as $x$ moves? That ``rate of change'' is the slope of the curve, and the derivative is the
tool that measures it.

The trick is to look at a tiny step. Move from $x$ to $x+h$. The output changes from $f(x)$ to
$f(x+h)$. The average rate of change over that step is the rise over the run,
\[
  \frac{f(x+h)-f(x)}{h}.
\]
This is the slope of the straight line through the two nearby points $(x,f(x))$ and $(x+h, f(x+h))$
--- a \emph{secant} line. Now shrink the step. As $h\to 0$ the secant pivots until it just kisses the
curve at $x$: it becomes the \emph{tangent} line, and its slope is the derivative,
\[
  f'(x) \;=\; \lim_{h\to 0}\frac{f(x+h)-f(x)}{h}.
\]
We also write it $\frac{\dd f}{\dd x}$, read ``$\dd f$ by $\dd x$,'' a notation that deliberately
looks like a ratio of tiny changes, because morally that is what it is.

\begin{intuitionbox}
The derivative answers ``if I nudge the input by a hair, how much does the output move, per unit of
nudge?'' If $f'(x)=3$, then near $x$ the function climbs three units of output for every one unit of
input. If $f'(x)=0$, the curve is momentarily flat --- a peak, a valley, or a plateau. If $f'(x)<0$,
the function is falling. Everything else in differential calculus is bookkeeping on top of this one
idea.
\end{intuitionbox}

A handful of derivatives cover almost everything we need in this book. The power rule says
$\frac{\dd}{\dd x}x^n = n x^{n-1}$; in particular constants have derivative $0$ and lines $ax$ have
derivative $a$. The exponential is its own derivative, $\frac{\dd}{\dd x}e^{x}=e^{x}$, which is the
secret reason the exponential and normal distributions are so analytically friendly. The logarithm
gives $\frac{\dd}{\dd x}\ln x = 1/x$. And two combining rules do the heavy lifting: the product rule
$(fg)'=f'g+fg'$ and the chain rule $\big(f(g(x))\big)' = f'(g(x))\,g'(x)$. The chain rule is the one
you will reach for most when transforming random variables --- see the change-of-variables chapter,
where the factor $\abs{g'}$ that appears in density transformations is the chain rule made visible.

\begin{figure}[H]
\centering
\begin{tikzpicture}
\begin{axis}[
  width=11cm, height=6.2cm,
  xlabel={$x$}, ylabel={$y$},
  domain=-0.2:3, samples=120,
  axis lines=left, ymin=-0.5, ymax=5,
  legend pos=north west, legend cell align=left,
  title={A function, a secant, and the tangent it becomes}
]
\addplot[densityred, very thick] {0.5*x*x} node[pos=0.92,above left,black] {$f(x)=\tfrac12 x^2$};
% tangent at x=1.5: f'=x=1.5, f=1.125, line y=1.5(x-1.5)+1.125
\addplot[sourceblue, thick, domain=0.3:2.7] {1.5*(x-1.5)+1.125};
\addlegendentry{tangent at $x=1.5$, slope $f'(1.5)=1.5$}
% secant from x=1 to x=2.5
\addplot[quotientorange, thick, dashed, domain=0.6:2.9] {((0.5*2.5*2.5)-(0.5*1*1))/(2.5-1)*(x-1)+0.5};
\addlegendentry{secant through $x=1$ and $x=2.5$}
\addplot[only marks, black, mark=*, mark size=1.4pt] coordinates {(1.5,1.125)};
\end{axis}
\end{tikzpicture}
\caption{The derivative is the limiting slope. The dashed secant connects two points on the curve; as
the points slide together it pivots into the solid tangent, whose slope is $f'(1.5)=1.5$.}
\label{chAcla:fig-tangent}
\end{figure}

\begin{connectionbox}
In probability the derivative shows up the moment a quantity is continuous. The probability density
function $f$ is literally the derivative of the cumulative distribution function $F$: $f(x)=F'(x)$.
That is the whole content of ``density'' --- it is the \emph{rate} at which probability accumulates as
$x$ increases. We unpack this in the continuous-random-variables chapter; here just notice that
``density = derivative of the CDF'' is the same tangent-line idea applied to a curve that happens to
measure probability. The companion \emph{Calculus From the Ground Up} treats limits and derivatives in
full.
\end{connectionbox}

\subsection{The integral as accumulated area}
\label{chAcla:integral}

Now turn the question around. Instead of ``how fast is $f$ changing,'' ask ``how much of $f$ has
piled up between $x=a$ and $x=b$?'' Geometrically: what is the area under the curve $y=f(x)$ from $a$
to $b$? That area is the definite integral
\[
  \int_a^b f(x)\,\dd x.
\]

The construction mirrors the derivative's limiting step, but in reverse spirit. Chop the interval
$[a,b]$ into many thin slices of width $\Delta x$. Over each slice the function is nearly constant, so
the area of that sliver is approximately height $\times$ width, $f(x_i)\,\Delta x$. Add up all the
slivers to get a \emph{Riemann sum}, $\sum_i f(x_i)\,\Delta x$, and then let the slices get
infinitely thin. In the limit the jagged staircase of rectangles becomes the smooth area under the
curve, and the summation sign $\sum$ relaxes into the integral sign $\int$ --- an elongated S, for
``sum.'' The $\dd x$ is the vanishing width left over from $\Delta x$.

\begin{figure}[H]
\centering
\begin{tikzpicture}
\begin{axis}[
  width=11cm, height=6cm,
  xlabel={$x$}, ylabel={$y$},
  domain=0:3, samples=120,
  axis lines=left, ymin=0, ymax=3.4,
  title={Rectangles approximate the area; the integral is their limit}
]
\addplot[densityred, very thick] {0.4*x*x+0.3} node[pos=0.95,above left,black]{$f(x)$};
\foreach \xa in {0,0.5,1,1.5,2,2.5} {
  \edef\temp{\noexpand\addplot[draw=sourceblue, fill=histfill, fill opacity=0.45]
    coordinates {(\xa,0) (\xa,{0.4*\xa*\xa+0.3}) (\xa+0.5,{0.4*\xa*\xa+0.3}) (\xa+0.5,0)} -- cycle;}
  \temp
}
\end{axis}
\end{tikzpicture}
\caption{A Riemann sum. Each rectangle has area (height)$\times$(width)$=f(x_i)\,\Delta x$. As the
widths shrink, the total rectangle area converges to $\int_0^3 f(x)\,\dd x$, the exact area under the
curve.}
\label{chAcla:fig-riemann}
\end{figure}

\begin{happeningbox}
An integral is a continuous sum. That single sentence is the bridge between the discrete and
continuous worlds of probability. When a random variable is discrete you compute its expectation by a
weighted sum, $\E[X]=\sum_x x\,\P(X=x)$. When it is continuous you compute the \emph{same} thing as a
weighted ``sum'' over a continuum, $\E[X]=\int x f(x)\,\dd x$. Nothing conceptual changed --- you are
still adding up ``value times how much probability sits there.'' Only the bookkeeping switched from
$\sum$ to $\int$ because there are now uncountably many values, each carrying an infinitesimal slice
of probability $f(x)\,\dd x$. The measure-theory chapters make this unification airtight; the integral
is the everyday face of it.
\end{happeningbox}

\subsection{The Fundamental Theorem of Calculus}
\label{chAcla:ftc}

The two halves of calculus --- slopes and areas --- look unrelated. The miracle, and it really is
one, is that they are inverse operations. This is the Fundamental Theorem of Calculus.

\begin{theorem}[Fundamental Theorem of Calculus]
\label{chAcla:thm-ftc}
Let $f$ be continuous on $[a,b]$. Define the ``area-so-far'' function
$\displaystyle A(x)=\int_a^x f(t)\,\dd t$. Then $A$ is differentiable and $A'(x)=f(x)$. Consequently,
if $F$ is any antiderivative of $f$ (so $F'=f$), then
\[
  \int_a^b f(x)\,\dd x \;=\; F(b)-F(a).
\]
\end{theorem}

\begin{whybox}
Why is the derivative of ``area so far'' equal to the height of the curve? Push the right edge of the
shaded region from $x$ to $x+h$. You add a thin strip of area. That strip is almost a rectangle of
width $h$ and height $f(x)$, so the added area is about $f(x)\cdot h$. Therefore
\[
  \frac{A(x+h)-A(x)}{h}\approx f(x),
\]
and as $h\to 0$ the approximation becomes exact: $A'(x)=f(x)$. The rate at which area accumulates is
exactly the current height of the curve --- which is just common sense once you see the strip.
Knowing this, evaluating an integral becomes a search for an antiderivative: find any $F$ with
$F'=f$, then subtract its values at the endpoints.
\end{whybox}

This is exactly the relationship between a density and a CDF. Because $f=F'$, the Fundamental Theorem
says $\P(a\le X\le b)=\int_a^b f(x)\,\dd x = F(b)-F(a)$. Differentiating a CDF gives the density;
integrating a density gives back probabilities. Calculus is the dictionary that translates between the
two descriptions of a continuous random variable.

\begin{example}[The exponential density integrates to one]
\label{chAcla:ex-expo}
The $\Expo(\lambda)$ density is $f(x)=\lambda e^{-\lambda x}$ for $x\ge 0$. An antiderivative of
$\lambda e^{-\lambda x}$ is $F(x)=-e^{-\lambda x}$, because by the chain rule $F'(x)=\lambda
e^{-\lambda x}$. Then the total probability is
\[
  \int_0^\infty \lambda e^{-\lambda x}\,\dd x
  = \Big[-e^{-\lambda x}\Big]_0^\infty
  = \big(0\big)-\big(-1\big) = 1,
\]
as a density must. The same antiderivative gives any tail: $\P(X>t)=\int_t^\infty \lambda e^{-\lambda
x}\,\dd x = e^{-\lambda t}$, the memoryless survival curve you meet in the exponential-distribution
section.
\end{example}

\subsection{Integration techniques you will actually meet}
\label{chAcla:techniques}

Two techniques recur often enough to name. The first is \emph{substitution}, which is the chain rule
read backwards. If you spot an inner function $u=g(x)$ whose derivative $g'(x)$ also appears (up to a
constant), substitute and the integral simplifies:
\[
  \int f(g(x))\,g'(x)\,\dd x = \int f(u)\,\dd u, \qquad u=g(x),\ \dd u = g'(x)\,\dd x.
\]
This is precisely the engine behind change-of-variables for densities: when $Y=g(X)$, the $\dd u =
g'(x)\,\dd x$ bookkeeping is where the Jacobian factor $\abs{g'}$ comes from.

The second is \emph{integration by parts}, the product rule read backwards:
\[
  \int u\,\dd v = uv - \int v\,\dd u.
\]
You use it to derive moments of distributions. For instance, $\E[X]$ for an exponential, or the
gamma-function identity $\Gamma(n)=(n-1)!$, both fall out of one or two integrations by parts.

\begin{example}[Mean of the exponential, by parts]
\label{chAcla:ex-mean-expo}
For $X\sim\Expo(\lambda)$,
\[
  \E[X]=\int_0^\infty x\,\lambda e^{-\lambda x}\,\dd x.
\]
Take $u=x$ (so $\dd u=\dd x$) and $\dd v=\lambda e^{-\lambda x}\,\dd x$ (so $v=-e^{-\lambda x}$).
Then
\[
  \E[X]=\Big[-x e^{-\lambda x}\Big]_0^\infty + \int_0^\infty e^{-\lambda x}\,\dd x
  = 0 + \Big[-\tfrac1\lambda e^{-\lambda x}\Big]_0^\infty = \frac1\lambda.
\]
The boundary term vanishes because $e^{-\lambda x}$ crushes the linear $x$ at infinity. So the mean
of an exponential is $1/\lambda$ --- a one-line computation once integration by parts is in hand.
\end{example}

\begin{historybox}
The integral sign $\int$ and the $\dd x$ notation are Leibniz's, introduced in the 1670s; the
elongated S stood for \emph{summa}, ``sum.'' Newton, working at the same time, thought of the same
ideas in terms of ``fluxions'' --- quantities flowing in time --- and the two were later embroiled in
a bitter priority dispute. For probability the relevant later hero is Henri Lebesgue, who in 1902
reinvented the integral by slicing the \emph{range} of $f$ into thin horizontal bands instead of the
domain into vertical strips. That seemingly small change is what lets us integrate against
probability measures and unify discrete and continuous expectation, and it is the foundation of the
measure-theoretic chapters near the end of this book.
\end{historybox}

\subsection{Improper integrals and why tails can be finite}
\label{chAcla:improper}

Probability densities live on infinite ranges --- the normal stretches over all of $\R$, the
exponential over $[0,\infty)$. Integrating to infinity is an \emph{improper} integral, defined as a
limit: $\int_a^\infty f = \lim_{b\to\infty}\int_a^b f$. The remarkable fact, which beginners find
counterintuitive, is that an infinitely long region can enclose a finite area, provided the curve
decays fast enough.

\begin{paradoxbox}
\textbf{Infinite length, finite area.} Consider $f(x)=1/x^2$ on $[1,\infty)$. The region under it
runs forever to the right, yet
\[
  \int_1^\infty \frac{1}{x^2}\,\dd x = \Big[-\tfrac1x\Big]_1^\infty = 0-(-1) = 1.
\]
A finite answer for an unbounded region! The resolution is that ``length'' and ``area'' are different
beasts: the strips get thin so fast that their areas form a convergent sum. But beware the other
direction. With $f(x)=1/x$ the same region has \emph{infinite} area, $\int_1^\infty \frac1x\,\dd x =
\lim_{b\to\infty}\ln b = \infty$, even though $1/x$ also decays to zero. The boundary between ``decays
fast enough'' and ``not quite'' is exactly $1/x$, and it is the same boundary that decides whether a
heavy-tailed distribution has a finite mean. A Cauchy distribution's density decays like $1/x^2$, so
it integrates to a valid probability of $1$; but $x$ times that density decays like $1/x$, so its
\emph{mean} integral diverges --- which is why the Cauchy has no mean at all. The St.\ Petersburg
paradox and heavy-tailed risk models in the inequalities chapter are this same arithmetic of
convergence.
\end{paradoxbox}

\subsection{Functions of several variables and multiple integrals}
\label{chAcla:multiple-integrals}

The moment two random variables share a story, we need functions of two inputs and integrals over
two-dimensional regions. A joint density $f(x,y)$ is a surface sitting above the plane; the
probability that $(X,Y)$ lands in a region $R$ is the \emph{volume} under that surface over $R$,
\[
  \P\big((X,Y)\in R\big) = \iint_R f(x,y)\,\dd x\,\dd y.
\]

You compute a double integral by integrating one variable at a time --- an \emph{iterated} integral.
Hold $y$ fixed, integrate over $x$ (sweeping out the area of a single slice), then integrate the
resulting function of $y$. Fubini's theorem guarantees that, for the well-behaved densities we use,
the order does not matter: $\iint = \int(\int \dd x)\dd y = \int(\int \dd y)\dd x$. This is the engine
behind \emph{marginalization}: to get the density of $X$ alone, integrate the joint density over all
of $y$, $f_X(x)=\int_{-\infty}^{\infty} f(x,y)\,\dd y$. You are summing the probability over every way
the other variable could have come out --- the law of total probability, written as an integral.

\begin{figure}[H]
\centering
\begin{tikzpicture}
\begin{axis}[
  width=9.5cm, height=6.5cm,
  xlabel={$x$}, ylabel={$y$}, zlabel={$f(x,y)$},
  domain=-2.2:2.2, y domain=-2.2:2.2, samples=24,
  view={40}{30}, colormap/viridis,
  title={A joint density as a surface; probability is volume underneath}
]
\addplot3[surf, opacity=0.85, faceted color=black!30]
  {0.159*exp(-0.5*(x*x+y*y))};
\end{axis}
\end{tikzpicture}
\caption{The standard bivariate normal density $f(x,y)=\frac{1}{2\pi}e^{-(x^2+y^2)/2}$. The
probability of any region $R$ in the $xy$-plane is the volume of the solid standing over $R$ and below
this bell-shaped surface. Marginalizing out $y$ flattens the surface onto the $x$-axis.}
\label{chAcla:fig-bivariate}
\end{figure}

\begin{connectionbox}
Polar coordinates earn their keep in probability through one famous trick: the integral
$\int_{-\infty}^{\infty} e^{-x^2/2}\,\dd x$ has no elementary antiderivative, yet its \emph{square}
becomes a double integral that switches to polar coordinates and evaluates cleanly to $2\pi$, giving
$\int e^{-x^2/2}\dd x = \sqrt{2\pi}$. That is exactly the normalizing constant of the standard normal,
and the substitution $r^2=x^2+y^2$, $\dd x\,\dd y = r\,\dd r\,\dd\theta$ is the reason the $\sqrt{2\pi}$
appears. The normal-distribution chapter does this computation; the multiple-integral machinery here
is what makes it legal.
\end{connectionbox}

\subsection{Gradients: derivatives in many directions}
\label{chAcla:gradient}

For a function of several variables, $g(x_1,\dots,x_n)$, there is no single ``slope'' --- the surface
tilts differently depending on which way you walk. The \emph{partial derivative} $\partial g/\partial
x_i$ measures the slope when you move along the $x_i$ axis only, holding the others fixed; it is just
an ordinary derivative with the other variables treated as constants. Collect all the partials into a
vector and you get the \emph{gradient},
\[
  \nabla g = \left(\frac{\partial g}{\partial x_1},\ \dots,\ \frac{\partial g}{\partial x_n}\right).
\]

\begin{intuitionbox}
The gradient points in the direction of steepest ascent, and its length is how steep that climb is.
Stand on a hillside; the gradient is the compass bearing of ``straight uphill,'' and its magnitude is
the grade. At a peak, valley, or saddle the ground is momentarily level in every direction, so the
gradient is the zero vector --- that is the multivariable version of ``set the derivative to zero to
find the maximum.''
\end{intuitionbox}

This is the workhorse of estimation. Maximum likelihood (the MLE chapter) finds the parameter that
maximizes a likelihood by setting its gradient to zero and solving. Modern machine learning (the ML
chapter) does the same in reverse: \emph{gradient descent} repeatedly steps in the direction of
$-\nabla g$, downhill, to minimize a loss function. Both are the one-variable idea ``optimum where
the derivative is zero'' lifted to many parameters at once. The companion \emph{Calculus From the
Ground Up} develops partial derivatives and gradients carefully; we only need the picture of ``uphill
arrow whose zero marks an optimum.''

\section{Part II: The Linear Algebra You Need}
\label{chAcla:la-intro}

If calculus is the mathematics of \emph{change}, linear algebra is the mathematics of \emph{many
quantities handled together}. The instant you study several random variables jointly --- their
spread, their correlations, how a Markov chain shuffles probability between states --- you are doing
linear algebra whether you call it that or not. This half builds the small toolkit the book relies on.

\subsection{Vectors and the dot product}
\label{chAcla:vectors}

A vector is an ordered list of numbers, $\mathbf{v}=(v_1,\dots,v_n)$, which you can picture as an
arrow from the origin to the point $(v_1,\dots,v_n)$ in $n$-dimensional space. In probability a vector
might hold a random sample $(X_1,\dots,X_n)$, a vector of means, or the probabilities of a chain's
states. Two operations make vectors useful: you can add them coordinatewise, and you can scale them by
a number. That is the entire definition of ``linear'' --- additions and scalings, nothing more exotic.

The single most important vector operation for us is the \emph{dot product},
\[
  \mathbf{u}\cdot\mathbf{v} = u_1 v_1 + u_2 v_2 + \cdots + u_n v_n = \sum_{i=1}^n u_i v_i.
\]
It eats two vectors and returns one number. Two facts give it meaning. First, a vector's length is
$\norm{\mathbf{v}} = \sqrt{\mathbf{v}\cdot\mathbf{v}}$ --- the Pythagorean theorem in $n$ dimensions.
Second, the dot product measures alignment: $\mathbf{u}\cdot\mathbf{v} =
\norm{\mathbf{u}}\,\norm{\mathbf{v}}\cos\theta$, where $\theta$ is the angle between the arrows. So a
dot product of zero means the vectors are perpendicular --- \emph{orthogonal} --- and a large positive
dot product means they point much the same way (Figure~\ref{chAcla:fig-dotproduct}).

\begin{happeningbox}
Correlation \emph{is} a dot product. Center two data vectors by subtracting their means, and the
sample correlation is exactly the cosine of the angle between the centered vectors:
\[
  \Corr(X,Y) = \frac{\sum_i (x_i-\bar x)(y_i-\bar y)}
                    {\sqrt{\sum_i (x_i-\bar x)^2}\,\sqrt{\sum_i (y_i-\bar y)^2}}
            = \cos\theta.
\]
Numerator: a dot product. Denominator: the two lengths. This is why correlation always lies in
$[-1,1]$ --- because a cosine does --- and why ``uncorrelated'' is the statistical name for
``orthogonal.'' The covariance-and-correlation chapter leans on this geometry; the whole notion of one
variable being ``perpendicular'' to another is the dot product talking.
\end{happeningbox}

\begin{figure}[H]
\centering
\begin{tikzpicture}[scale=1.0, >=Stealth]
\draw[->, very thick, sourceblue] (0,0) -- (3.4,0.6) node[right]{$\mathbf{u}$};
\draw[->, very thick, targetgreen] (0,0) -- (1.8,2.6) node[above]{$\mathbf{v}$};
\draw[densityred, thick] (0.7,0.123) arc[start angle=10, end angle=55, radius=0.71];
\node[densityred] at (1.05,0.62) {$\theta$};
\node[black, align=left] at (4.7,1.7)
  {$\mathbf{u}\cdot\mathbf{v}=\norm{\mathbf{u}}\norm{\mathbf{v}}\cos\theta$};
\end{tikzpicture}
\caption{The dot product as alignment. Two vectors meet at angle $\theta$; their dot product is the
product of their lengths times $\cos\theta$. When the centered data vectors are used, this $\cos\theta$
is exactly the sample correlation, which is why correlation is pinned to $[-1,1]$ and why
``uncorrelated'' means ``perpendicular.''}
\label{chAcla:fig-dotproduct}
\end{figure}

\subsection{Matrices as linear transformations}
\label{chAcla:matrices}

A matrix is a rectangular grid of numbers, but the right way to think about it is as a \emph{machine}:
feed it a vector, it returns another vector. A matrix $A$ acts on $\mathbf{x}$ by the rule that the
$i$-th output coordinate is the dot product of the $i$-th row of $A$ with $\mathbf{x}$. Matrix
multiplication is just the composition of two such machines: $AB$ means ``do $B$, then $A$.''

The crucial property is \emph{linearity}: $A(\mathbf{x}+\mathbf{y})=A\mathbf{x}+A\mathbf{y}$ and
$A(c\mathbf{x})=cA\mathbf{x}$. Geometrically, a matrix sends lines to lines and keeps the origin
fixed; it can rotate, stretch, shear, and reflect space, but it cannot bend it. This is exactly the
linearity that makes $\E$ behave well: expectation of a linear combination is the linear combination
of expectations, $\E[A\mathbf{X}] = A\,\E[\mathbf{X}]$, treating the random vector $\mathbf{X}$ just
like any vector the matrix acts on.

\begin{example}[A $2\times 2$ map in action]
\label{chAcla:ex-matrix-action}
Let $A=\begin{pmatrix} 2 & 0 \\ 0 & \tfrac12 \end{pmatrix}$. Acting on $\mathbf{x}=(x_1,x_2)$ it
returns $(2x_1,\ \tfrac12 x_2)$: it stretches the horizontal direction by $2$ and squashes the
vertical by half. A circle of data points becomes an ellipse twice as wide and half as tall. Keep
this picture --- ``a matrix turns a circle into an ellipse'' --- because it is precisely how a
covariance matrix reshapes a round cloud of independent noise into the tilted elliptical cloud of
correlated data.
\end{example}

\subsection{Transition matrices and the stochastic special case}
\label{chAcla:transition}

A Markov chain is a matrix in disguise. List the states $1,\dots,n$ and let $P_{ij}=\P(\text{go to }j
\given \text{now in }i)$. The matrix $P$ is \emph{stochastic}: every entry is a probability in
$[0,1]$ and every row sums to $1$. If $\boldsymbol\pi$ is a row vector holding the current
distribution over states, then one step of the chain is one matrix multiplication, $\boldsymbol\pi
\mapsto \boldsymbol\pi P$, and $k$ steps is $\boldsymbol\pi P^k$. The entire dynamics of the chain ---
where probability flows, where it pools --- is encoded in powers of one matrix.

\begin{figure}[H]
\centering
\begin{tikzpicture}[>=Stealth, node distance=3.2cm, every node/.style={font=\small}]
\node[circle, draw=sourceblue, very thick, fill=blue!5, minimum size=1.2cm] (S) {Sun};
\node[circle, draw=quotientorange, very thick, fill=orange!7, minimum size=1.2cm, right=of S] (R) {Rain};
\draw[->, thick] (S) to[bend left=18] node[above]{$0.3$} (R);
\draw[->, thick] (R) to[bend left=18] node[below]{$0.5$} (S);
\draw[->, thick] (S) to[out=120,in=60,looseness=5] node[above]{$0.7$} (S);
\draw[->, thick] (R) to[out=120,in=60,looseness=5] node[above]{$0.5$} (R);
\end{tikzpicture}
\caption{A two-state weather chain and its transition matrix $P=\left(\begin{smallmatrix} 0.7 & 0.3
\\ 0.5 & 0.5 \end{smallmatrix}\right)$. Each row sums to $1$. Multiplying a distribution vector by $P$
advances the weather one day; the steady state is the vector $\boldsymbol\pi$ fixed by $\boldsymbol\pi
P=\boldsymbol\pi$.}
\label{chAcla:fig-markov}
\end{figure}

That steady-state condition, $\boldsymbol\pi P = \boldsymbol\pi$, says $\boldsymbol\pi$ is unchanged by
the matrix --- it is a special vector for $P$. Naming such vectors is the next idea.

\subsection{Eigenvectors and eigenvalues}
\label{chAcla:eigen}

Most vectors get knocked off their direction when a matrix acts on them. But for any square matrix
there are usually a few special directions that survive: vectors $\mathbf{v}$ that the matrix merely
stretches, without rotating. Such a vector satisfies
\[
  A\mathbf{v} = \lambda\mathbf{v}
\]
for some number $\lambda$. We call $\mathbf{v}$ an \emph{eigenvector} and $\lambda$ its
\emph{eigenvalue} (German \emph{eigen}, ``own'' or ``characteristic''). The eigenvector names a
direction the matrix leaves alone; the eigenvalue says by how much it scales along that direction.

\begin{intuitionbox}
Picture the matrix as a flow that pushes every arrow around. Eigenvectors are the directions where the
push is purely outward or inward --- no swirl, just a clean stretch by the factor $\lambda$. If
$\abs{\lambda}>1$ that direction grows under repeated application; if $\abs{\lambda}<1$ it shrinks; if
$\lambda=1$ it is perfectly preserved. Repeatedly applying the matrix amplifies the largest-eigenvalue
direction and washes the rest out, which is why eigenvectors govern long-run behavior.
\end{intuitionbox}

This is the secret behind Markov-chain steady states. A stochastic matrix always has eigenvalue
$\lambda=1$, and its eigenvector is exactly the stationary distribution $\boldsymbol\pi$ with
$\boldsymbol\pi P=\boldsymbol\pi$ --- a $\lambda=1$ direction that the chain preserves forever. All
other eigenvalues have $\abs{\lambda}<1$, so their contributions decay; that decay is precisely why a
well-behaved chain forgets its starting point and converges to $\boldsymbol\pi$. The
steady-state chapter proves this; the eigenvalue picture is why it is true.

\begin{aestheticbox}
Eigenvectors are how a matrix tells you its own preferred coordinate system. In the right basis ---
the eigenvector basis --- a complicated transformation becomes nothing but independent stretches along
each axis, a diagonal matrix. Half of applied linear algebra is the art of finding that basis, because
in it the hard problem falls apart into easy one-dimensional pieces. We are about to see this with
covariance matrices, where the eigenvectors are the axes of the data cloud's ellipse.
\end{aestheticbox}

\subsection{Symmetric and positive-definite matrices}
\label{chAcla:posdef}

The matrices that appear in probability are not arbitrary; they are almost always \emph{symmetric},
meaning $A=A^\top$ (the entry in row $i$ column $j$ equals the entry in row $j$ column $i$). A
covariance matrix is the prime example: its $(i,j)$ entry is $\Cov(X_i,X_j)=\Cov(X_j,X_i)$, manifestly
symmetric. Symmetry is a gift, because it forces excellent behavior.

\begin{theorem}[Spectral theorem, the version we use]
\label{chAcla:thm-spectral}
Every real symmetric matrix $A$ has only real eigenvalues, and its eigenvectors can be chosen mutually
orthogonal and of unit length. Collecting them as the columns of an orthogonal matrix $Q$ (so
$Q^\top Q = I$) and the eigenvalues on the diagonal of $\Lambda$, we can write
\[
  A = Q\,\Lambda\,Q^\top.
\]
\end{theorem}

This factorization says any symmetric matrix is ``rotate into the eigenvector frame ($Q^\top$), stretch
each axis by its eigenvalue ($\Lambda$), rotate back ($Q$).'' It is the rigorous form of ``a matrix
turns a circle into an ellipse,'' and it is exactly the decomposition $\Sigma = Q\Lambda Q^\top$ that
the multivariate-normal chapter uses to draw the elliptical contours of a Gaussian and to sample from
it.

A symmetric matrix is \emph{positive-definite} if all its eigenvalues are strictly positive, or
equivalently if $\mathbf{x}^\top A\,\mathbf{x} > 0$ for every nonzero vector $\mathbf{x}$ (it is
\emph{positive-semidefinite} if we allow $\ge 0$). The quantity $\mathbf{x}^\top A\,\mathbf{x}$ is a
\emph{quadratic form} --- the matrix generalization of $a x^2$ --- and ``positive-definite'' is the
matrix way of saying ``opens upward like a bowl.''

\begin{whybox}
Why must a covariance matrix be positive-semidefinite? Take any vector of weights $\mathbf{a}$ and
form the scalar random variable $Y=\mathbf{a}^\top \mathbf{X}=\sum_i a_i X_i$. Its variance cannot be
negative --- variance is an average of squared deviations. But a direct expansion gives
\[
  0 \le \Var(Y) = \Var\!\Big(\sum_i a_i X_i\Big)
      = \sum_{i,j} a_i a_j \Cov(X_i,X_j)
      = \mathbf{a}^\top \Sigma\,\mathbf{a}.
\]
So $\mathbf{a}^\top \Sigma\,\mathbf{a}\ge 0$ for \emph{every} $\mathbf{a}$ --- which is the definition
of positive-semidefinite. The abstract algebraic condition is nothing but ``no combination of your
variables can have negative variance.'' That is also why you cannot hand a statistics package an
arbitrary symmetric matrix and call it a covariance matrix: if it has a negative eigenvalue, some
linear combination would have negative variance, which is impossible.
\end{whybox}

\begin{figure}[H]
\centering
\begin{tikzpicture}
\begin{axis}[
  width=8.4cm, height=6.2cm,
  xlabel={$x_1$}, ylabel={$x_2$},
  axis equal image, axis lines=middle,
  xmin=-3.4, xmax=3.4, ymin=-3.4, ymax=3.4,
  xtick=\empty, ytick=\empty,
  title={Eigenvectors are the axes of the covariance ellipse}
]
% ellipse for Sigma with eigenvalues 4 (along (1,1)/sqrt2) and 1 (along (-1,1)/sqrt2)
% point = 2 cos(t) v1 + 1 sin(t) v2, v1=(1,1)/sqrt2, v2=(-1,1)/sqrt2
\addplot[densityred, very thick, domain=0:360, samples=120]
  ({2*cos(x)*0.7071 + 1*sin(x)*(-0.7071)}, {2*cos(x)*0.7071 + 1*sin(x)*0.7071});
% major eigen-axis (1,1)
\addplot[sourceblue, very thick, -{Stealth}] coordinates {(0,0) (1.414,1.414)};
\node[sourceblue] at (axis cs:1.9,1.3) {$\mathbf{v}_1$};
% minor eigen-axis (-1,1)
\addplot[targetgreen, very thick, -{Stealth}] coordinates {(0,0) (-0.707,0.707)};
\node[targetgreen] at (axis cs:-1.25,0.75) {$\mathbf{v}_2$};
\end{axis}
\end{tikzpicture}
\caption{A correlated Gaussian's contour. The eigenvectors $\mathbf{v}_1,\mathbf{v}_2$ of the
covariance matrix point along the ellipse's principal axes; the square roots of the eigenvalues give
the axis half-lengths. This is the spectral theorem made visible, and it is exactly what principal
component analysis extracts from data.}
\label{chAcla:fig-ellipse}
\end{figure}

\begin{applicationbox}
\textbf{Principal Component Analysis (PCA).} Given a cloud of data, compute its covariance matrix
$\Sigma$ and find the eigenvectors. The eigenvector with the largest eigenvalue points along the
direction of greatest variance --- the long axis of the data ellipse --- and the eigenvalue \emph{is}
the variance captured along it. Keeping the top few eigenvectors compresses high-dimensional data to a
few coordinates while preserving as much variance as possible. PCA is, almost in its entirety, ``apply
the spectral theorem to a covariance matrix,'' which is why the machine-learning chapter can treat it
in a couple of pages: all the heavy lifting was done the moment we knew $\Sigma=Q\Lambda Q^\top$.
\end{applicationbox}

\begin{historybox}
The word \emph{eigenvalue} is a half-translation of the German \emph{Eigenwert} (``own value''),
which entered mathematics through David Hilbert's work on integral equations around 1904. The deeper
story is older: the ``principal axes'' of a quadratic form were studied by Cauchy in the 1820s and by
Lagrange before him, in the guise of finding the axes of an ellipse or the principal stresses in a
solid. Karl Pearson, one of the founders of mathematical statistics, reinvented the same axes in 1901
under the name ``principal components'' --- the very PCA above --- without realizing he was rediscovering
the spectral theorem. Probability, mechanics, and pure algebra all converged on the same eigenvectors.
\end{historybox}

\subsection{Orthogonal projection}
\label{chAcla:projection}

The last tool is projection: given a vector $\mathbf{y}$ and a subspace (say, the line or plane
spanned by some other vectors), find the point in the subspace closest to $\mathbf{y}$. Geometrically
you drop a perpendicular from $\mathbf{y}$ onto the subspace; the foot of that perpendicular is the
projection $\hat{\mathbf{y}}$, and the leftover $\mathbf{y}-\hat{\mathbf{y}}$ is orthogonal to
everything in the subspace.

For projection onto the line through a unit vector $\mathbf{u}$, the formula is just a dot product
times a direction,
\[
  \hat{\mathbf{y}} = (\mathbf{y}\cdot\mathbf{u})\,\mathbf{u}.
\]
The number $\mathbf{y}\cdot\mathbf{u}$ is ``how much of $\mathbf{y}$ lies along $\mathbf{u}$,'' and we
lay that much off in the $\mathbf{u}$ direction. The key property is that $\hat{\mathbf{y}}$ is the
\emph{closest} point in the subspace --- projection minimizes squared distance.

\begin{happeningbox}
Least squares is a projection, and so is conditional expectation. When you fit a line by least
squares (the estimation chapter), you are projecting the vector of observed responses onto the subspace
of fitted values, choosing the prediction that minimizes the sum of squared errors --- exactly the
closest point. More strikingly, the conditional expectation $\E[Y\given X]$ is the orthogonal
projection of $Y$ onto the space of functions of $X$: among all predictors built from $X$, it is the
one with the smallest mean squared error. ``Best prediction'' and ``perpendicular drop'' are the same
statement in different languages. The residual $Y-\E[Y\given X]$ being uncorrelated with every
function of $X$ is just ``the leftover is orthogonal to the subspace,'' which the conditional-
expectation chapter calls the orthogonality principle.
\end{happeningbox}

\begin{figure}[H]
\centering
\begin{tikzpicture}[scale=1.0, >=Stealth]
\draw[->, thick, gray] (-0.4,0) -- (5.2,0) node[right]{subspace (span of $\mathbf{u}$)};
\draw[->, very thick, sourceblue] (0,0) -- (3.2,2.4) node[above right]{$\mathbf{y}$};
\draw[->, very thick, targetgreen] (0,0) -- (3.2,0) node[below]{$\hat{\mathbf{y}}$};
\draw[->, very thick, densityred] (3.2,0) -- (3.2,2.4)
  node[midway, right]{$\mathbf{y}-\hat{\mathbf{y}}$ (residual)};
\draw (2.95,0) -- (2.95,0.25) -- (3.2,0.25); % right-angle mark
\end{tikzpicture}
\caption{Orthogonal projection of $\mathbf{y}$ onto a subspace. The projection $\hat{\mathbf{y}}$ is
the closest point; the residual $\mathbf{y}-\hat{\mathbf{y}}$ meets the subspace at a right angle.
Least-squares regression and conditional expectation are both this picture.}
\label{chAcla:fig-projection}
\end{figure}

\begin{connectionbox}
The Pythagorean theorem returns one last time, now as a probabilistic identity. Because the residual
is orthogonal to the prediction, $\norm{\mathbf{y}}^2 = \norm{\hat{\mathbf{y}}}^2 +
\norm{\mathbf{y}-\hat{\mathbf{y}}}^2$. Translated into variances, this is the law of total variance,
$\Var(Y) = \Var\big(\E[Y\given X]\big) + \E\big[\Var(Y\given X)\big]$ --- ``explained plus
unexplained.'' Geometry and probability are telling the same story: total spread splits into the part
captured by the prediction and the part left in the residual, and they are perpendicular, so their
sizes add in squares. The conditional-expectation chapter develops this; \emph{Linear Algebra From the
Ground Up} develops projections and orthogonality in full generality.
\end{connectionbox}

\section{How the Pieces Fit the Book}
\label{chAcla:summary}

Step back and the two halves of this appendix line up neatly with the two faces of probability.

\begin{focusbox}
\textbf{Calculus, where it lands.}
\begin{itemize}
\item \emph{Derivative} $=$ density from CDF, $f=F'$; the Jacobian $\abs{g'}$ in change of variables.
\item \emph{Integral} $=$ continuous expectation $\int x f(x)\,\dd x$ and probabilities as areas
  $\int_a^b f$.
\item \emph{Fundamental Theorem} $=$ the two-way street between density and CDF.
\item \emph{Multiple integrals} $=$ joint densities, marginalization, the $\sqrt{2\pi}$ of the normal.
\item \emph{Gradient} $=$ maximum likelihood (set $\nabla=0$) and gradient descent in machine learning.
\end{itemize}

\textbf{Linear algebra, where it lands.}
\begin{itemize}
\item \emph{Dot product} $=$ correlation as a cosine; uncorrelated means orthogonal.
\item \emph{Matrix} $=$ linear maps, $\E[A\mathbf{X}]=A\E[\mathbf{X}]$, Markov transition steps.
\item \emph{Eigenvectors} $=$ Markov steady states ($\lambda=1$) and long-run behavior.
\item \emph{Spectral theorem / positive-definite} $=$ covariance matrices, Gaussian ellipses, PCA.
\item \emph{Projection} $=$ least squares and conditional expectation as best prediction.
\end{itemize}
\end{focusbox}

None of these tools is exotic, and you need only the working picture of each: a derivative is a slope,
an integral is an area, a gradient is an uphill arrow, a dot product measures alignment, a matrix
reshapes space, an eigenvector is a direction left alone, and a projection is the closest point. Carry
those seven images and the calculus and linear algebra in this book will read as common sense dressed
in notation. When you want the rigorous development behind any of them, the companion \emph{Calculus
From the Ground Up} and \emph{Linear Algebra From the Ground Up} are written in the same spirit as this
one.

\section{Exercises}
\label{chAcla:exercises}

\begin{exercise}[Density is the derivative of the CDF]
\label{chAcla:exr-density}
The CDF of an $\Expo(\lambda)$ random variable is $F(x)=1-e^{-\lambda x}$ for $x\ge 0$. Differentiate
$F$ to recover the density $f$, and confirm it equals $\lambda e^{-\lambda x}$. Then explain in one
sentence why ``density is the derivative of the CDF'' is just the Fundamental Theorem of Calculus.
\end{exercise}

\begin{exercise}[A density that integrates to one]
\label{chAcla:exr-normalize}
A random variable has density $f(x)=c\,x^2$ on $[0,2]$ and $0$ elsewhere. Find the constant $c$ by
requiring $\int_0^2 f(x)\,\dd x = 1$. Then compute $\E[X]=\int_0^2 x f(x)\,\dd x$.
\end{exercise}

\begin{exercise}[Integration by parts for a moment]
\label{chAcla:exr-parts}
Use integration by parts to evaluate $\int_0^\infty x^2\,\lambda e^{-\lambda x}\,\dd x$, the second
moment of an $\Expo(\lambda)$. (Hint: take $u=x^2$, integrate the rest; you will need to integrate by
parts a second time, or reuse the mean from Example~\ref{chAcla:ex-mean-expo}.) Confirm your answer is
$2/\lambda^2$, and deduce that $\Var(X)=1/\lambda^2$.
\end{exercise}

\begin{exercise}[Convergence of a tail integral]
\label{chAcla:exr-tail}
For which exponents $p>0$ does $\int_1^\infty x^{-p}\,\dd x$ converge to a finite value, and what is
that value when it does? Relate your finding to the paradox box on infinite length but finite area,
and explain why a density decaying like $1/x^2$ can be a valid density yet have an infinite mean.
\end{exercise}

\begin{exercise}[Marginalizing a joint density]
\label{chAcla:exr-marginal}
Let $f(x,y)=4xy$ on the unit square $0\le x\le 1$, $0\le y\le 1$, and $0$ elsewhere. Verify it
integrates to $1$ over the square. Then find the marginal density of $X$ by computing
$f_X(x)=\int_0^1 f(x,y)\,\dd y$, and check that the marginal itself integrates to $1$.
\end{exercise}

\begin{exercise}[A gradient set to zero]
\label{chAcla:exr-gradient}
Consider $g(a,b)=\sum_{i=1}^n (y_i - a - b x_i)^2$, the sum of squared errors of a line $a+bx$ fit to
data $(x_i,y_i)$. Compute the partial derivatives $\partial g/\partial a$ and $\partial g/\partial b$,
set both to zero, and show that the resulting equations are the normal equations of least squares.
Explain how this is the gradient-equals-zero idea from Section~\ref{chAcla:gradient}.
\end{exercise}

\begin{exercise}[Correlation as a cosine]
\label{chAcla:exr-cosine}
Take the two centered vectors $\mathbf{u}=(1,-1,0)$ and $\mathbf{v}=(2,0,-2)$. Compute
$\mathbf{u}\cdot\mathbf{v}$, the two lengths, and hence $\cos\theta$. Interpret the result as a
sample correlation. What would $\mathbf{u}\cdot\mathbf{v}=0$ have meant about the two variables?
\end{exercise}

\begin{exercise}[Eigenvector of a transition matrix]
\label{chAcla:exr-eigen}
For the weather chain $P=\left(\begin{smallmatrix} 0.7 & 0.3 \\ 0.5 & 0.5 \end{smallmatrix}\right)$
of Figure~\ref{chAcla:fig-markov}, find the row vector $\boldsymbol\pi=(\pi_1,\pi_2)$ with
$\boldsymbol\pi P=\boldsymbol\pi$ and $\pi_1+\pi_2=1$. Confirm $\boldsymbol\pi$ is a left eigenvector
of $P$ with eigenvalue $1$, and interpret it as the long-run fraction of sunny days.
\end{exercise}

\begin{exercise}[Positive-definiteness, by hand]
\label{chAcla:exr-posdef}
Let $A=\begin{pmatrix} 2 & 1 \\ 1 & 2 \end{pmatrix}$. Compute the quadratic form
$\mathbf{x}^\top A\,\mathbf{x}$ for $\mathbf{x}=(x_1,x_2)$, write it as a sum that is visibly
nonnegative, and conclude $A$ is positive-definite. Then find the eigenvalues of $A$ and verify they
are both positive, consistent with the criterion in Section~\ref{chAcla:posdef}.
\end{exercise}

\begin{exercise}[Projection minimizes distance (proof)]
\label{chAcla:exr-projection-proof}
Let $\mathbf{u}$ be a unit vector and $\hat{\mathbf{y}}=(\mathbf{y}\cdot\mathbf{u})\mathbf{u}$ the
projection of $\mathbf{y}$ onto the line through $\mathbf{u}$. For any scalar $t$, expand
$\norm{\mathbf{y}-t\mathbf{u}}^2 = \norm{\mathbf{y}}^2 - 2t(\mathbf{y}\cdot\mathbf{u}) + t^2$ and show
it is minimized at $t=\mathbf{y}\cdot\mathbf{u}$. Conclude that $\hat{\mathbf{y}}$ is the closest point
on the line, and that the residual $\mathbf{y}-\hat{\mathbf{y}}$ is orthogonal to $\mathbf{u}$.
Explain how this is the geometric heart of least-squares regression.
\end{exercise}
